% 附录 B 电磁学（Electromagnetism）
\chapter{电磁学}

本附录简要回顾电磁学中的若干关键结果。这些结果的详细推导请查阅延伸阅读所列
文献。

\section{磁矩}

磁通密度为 $\mathbf{B}$ 的磁场中，载流 I、长 $\mathrm{d}\mathbf{r}$ 的直导线段
所受洛伦兹力 $\mathrm{d}\mathbf{F}$ 为 $\mathrm{d}\mathbf{F}=I\,\mathrm{d}\mathbf{r}\times\mathbf{B}$
（见图 B.1）。

\begin{figure}
\centering
\includegraphics[width=0.42\textwidth]{figB_01.pdf}
\caption{电流元所受的力为 $\mathrm{d}\mathbf{F}=I\,\mathrm{d}\mathbf{r}\times\mathbf{B}$。}
\end{figure}

\begin{figure}
\centering
\includegraphics[width=0.42\textwidth]{figB_02.pdf}
\caption{取向垂直于 z 方向的正方形元电流环（边长 $\mathrm{d}x$ 与 $\mathrm{d}y$）。}
\end{figure}

这一结果可用于计算图 B.2 中元电流环所受的力偶：
\begin{equation*}
\begin{aligned}
\mathrm{d}G_x &= -I\,\mathrm{d}x\,\mathrm{d}y\,B_y\\
\mathrm{d}G_y &= I\,\mathrm{d}x\,\mathrm{d}y\,B_x,
\end{aligned}\tag{B.1}
\end{equation*}
于是定义环的磁矩 $\mathrm{d}\mu_z=I\,\mathrm{d}x\,\mathrm{d}y$，并对环的任意取向
重复该论证，力偶 $\mathrm{d}\mathbf{G}$ 为
\begin{equation*}
\mathrm{d}\mathbf{G} = \mathrm{d}\boldsymbol{\mu}\times\mathbf{B}.\tag{B.2}
\end{equation*}
对有限大小的环，积分得
\begin{equation*}
\mathbf{G} = \boldsymbol{\mu}\times\mathbf{B}.\tag{B.3}
\end{equation*}
只有磁场非均匀时环上才有净力。图 B.2 中元环在 x 方向所受的力为
\begin{equation*}
\mathrm{d}F_x = I\,\mathrm{d}y\left(B_z+\frac{\partial B_z}{\partial x}\mathrm{d}x\right) - I\,\mathrm{d}y\,B_z = \mathrm{d}\mu_z\frac{\partial B_z}{\partial x},\tag{B.4}
\end{equation*}
故
\begin{equation*}
\mathrm{d}\mathbf{F} = (\mathrm{d}\boldsymbol{\mu}\cdot\nabla)\mathbf{B}\tag{B.5}
\end{equation*}
对有限大小的磁矩积分得
\begin{equation*}
\mathbf{F} = (\boldsymbol{\mu}\cdot\nabla)\mathbf{B}.\tag{B.6}
\end{equation*}
因此只有非均匀磁场才能对磁矩施加力。

与磁场 B 排齐的磁矩 $\boldsymbol{\mu}$ 所受力矩为零（由式 B.3，$\boldsymbol{\mu}$
与 B 平行时 G=0）。因此只能通过施加力矩使 $\boldsymbol{\mu}$ 相对 B 转开。力矩
把 $\boldsymbol{\mu}$ 转到与磁场成 $\theta$ 角所做的功 W 为
\begin{equation*}
W = \int_{0}^{\theta}G\,\mathrm{d}\theta' = \int_{0}^{\theta}\mu B\sin\theta'\,\mathrm{d}\theta' = \mu B(1-\cos\theta).\tag{B.7}
\end{equation*}
磁矩的势能 U 因此可写为 $-\mu B\cos\theta$（略去常数项）。故可写
\begin{equation*}
U = -\boldsymbol{\mu}\cdot\mathbf{B}.\tag{B.8}
\end{equation*}
非均匀场 B 中磁矩 $\boldsymbol{\mu}$ 沿 x 方向所受的力按式 B.6 为
$\mu\,\partial B/\partial x$。要阻止它沿 x 运动，须施加恢复力
$-\mu\,\partial B/\partial x$。若允许它移动 $\delta x$ 进入更强的场区，该力做的功为
\begin{equation*}
\mathrm{d}W = -\mu\frac{\partial B}{\partial x}\delta x = -\mu\,\delta B,\tag{B.9}
\end{equation*}
这是样品自由能的增量（与习题 2.10 的结果一致）。

原点处磁矩产生的磁场为
\begin{equation*}
\mathbf{B}(\mathbf{r}) = -\mu_0\nabla\left(\frac{\boldsymbol{\mu}\cdot\mathbf{r}}{4\pi r^{3}}\right)\tag{B.10}
\end{equation*}
磁矢势（取最方便的规范）为
\begin{equation*}
\mathbf{A}(\mathbf{r}) = \frac{\mu_0}{4\pi}\frac{\boldsymbol{\mu}\times\mathbf{r}}{r^{3}}.\tag{B.11}
\end{equation*}

\section{自由空间中的麦克斯韦方程组}

自由空间中的麦克斯韦方程组为
\begin{align*}
\nabla\cdot\boldsymbol{\varepsilon} &= \rho/\epsilon_0\tag{B.12}\\
\nabla\cdot\mathbf{B} &= 0\tag{B.13}\\
\nabla\times\boldsymbol{\varepsilon} &= -\frac{\partial\mathbf{B}}{\partial t}\tag{B.14}\\
\nabla\times\mathbf{B} &= \mu_0\mathbf{J} + \epsilon_0\mu_0\frac{\partial\boldsymbol{\varepsilon}}{\partial t},\tag{B.15}
\end{align*}
描述电场 $\boldsymbol{\varepsilon}$、磁感应强度 B、电荷密度 $\rho$ 与电流密度
$\mathbf{J}$ 之间的相互关系。式 B.12 表明电场从正电荷发散、向负电荷汇聚；电荷
密度因此是电场的源或汇（见图 B.3）。式 B.13 表明磁场没有这种发散：不存在磁荷
（磁单极），B 场线只能构成回路，永不在任何地方起始或终止。式 B.14 表明：只有
在磁场随时间变化的空间区域周围才有电场回路——这引出电磁感应的法拉第定律。
不存在变化电场时，式 B.15 表明电流周围存在磁感应的回路。注意本节一切都是对
自由空间中的电磁场而言。
\bio{James Clerk Maxwell\\（1831--1879）}

\begin{figure}
\centering
\includegraphics[width=0.42\textwidth]{figB_03.pdf}
\caption{电场从正电荷发散、向负电荷汇聚。}
\end{figure}

\section{自由电流与束缚电流}

把物质纳入图景之后，事情变得稍微复杂，值得弄清原委。这里只讨论对磁性的效应，
尽管相应的电场推导背后是类似的思路。真实物质中含有一种与我们熟悉的、沿导线流动
的电流不同类型的电流：原子中含有电子绕核运动形成的微观束缚电流（bound
currents）。于是式 B.15 中的总电流 J 可以分解为两个分量
\begin{equation*}
\mathbf{J} = \mathbf{J}_{\text{free}} + \mathbf{J}_{\text{bound}},\tag{B.16}
\end{equation*}
分别来自自由电流（沿铜线流动的 10 安培）与束缚电流（绕原子旋转的电子）。

若 M 在样品中均匀，流过的唯一净束缚电流绕材料边缘。若 M 非均匀，材料内部会
产生束缚电流。我们需要束缚电流与磁化强度之间的一般关系，即
\begin{equation*}
\mathbf{J}_{\text{bound}} = \nabla\times\mathbf{M}.\tag{B.17}
\end{equation*}
现在概述式 B.17 的推导（更多细节见延伸阅读）。原点处单个点偶极子的矢势由式 B.11
给出。考虑体积 V 内的磁化样品，位置 $\mathbf{r}'$ 处磁化强度为
$\mathbf{M}(\mathbf{r}')$。位置 $\mathbf{r}$ 处的矢势为
\begin{equation*}
\mathbf{A}(\mathbf{r}) = \frac{\mu_0}{4\pi}\int_V\mathrm{d}^{3}r'\,\frac{\mathbf{M}(\mathbf{r}')\times(\mathbf{r}-\mathbf{r}')}{|\mathbf{r}-\mathbf{r}'|^{3}}.\tag{B.18}
\end{equation*}
利用 $\nabla_r(1/r)=-\mathbf{r}/r^{3}$、矢量恒等式
\begin{equation*}
\nabla\times(f\mathbf{g}) = f\nabla\times\mathbf{g} - \mathbf{g}\times\nabla f,\tag{B.19}
\end{equation*}
（f 与 g 分别是标量与矢量函数）并把对 $\mathbf{r}'$ 的积分变换为对 $\mathbf{r}$
的积分，此方程可以简化。结果是位置 $\mathbf{r}$ 处的矢势可写为
\begin{align*}
\mathbf{A}(\mathbf{r}) &= \frac{\mu_0}{4\pi}\int_V\frac{\nabla\times\mathbf{M}\,\mathrm{d}\tau}{r} + \frac{\mu_0}{4\pi}\int_S\nabla\times\left(\frac{\mathbf{M}}{r}\right)\mathrm{d}\tau\tag{B.20}\\
&= \frac{\mu_0}{4\pi}\int_V\frac{\nabla\times\mathbf{M}\,\mathrm{d}\tau}{r} + \frac{\mu_0}{4\pi}\int_V\frac{\mathbf{M}\times\mathrm{d}\mathbf{S}}{r}\tag{B.21}\\
&= \frac{\mu_0}{4\pi}\int_V\frac{\mathbf{J}_{\text{bound}}\,\mathrm{d}\tau}{r} + \frac{\mu_0}{4\pi}\int_S\frac{\mathbf{K}_{\text{bound}}\,\mathrm{d}S}{r}\tag{B.22}
\end{align*}
第一个积分遍及样品体积，第二个积分遍及样品表面。式 B.22 中的束缚电流与式 B.17
一致。面电流为 $\mathbf{K}_{\text{bound}}=\mathbf{M}\times\hat{\mathbf{n}}$，
矢量 $\hat{\mathbf{n}}$ 垂直于表面。

式 B.15 的问题在于：B 的旋度不仅联系自由电流密度，还联系更棘手的束缚电流密度。
这促使我们定义一个只“看见”自由电流密度的新“修正”场——尽管如下所示这只是
部分成功。利用束缚体电流的结果，定义新场 H（称为磁场强度，或干脆叫 H 场）：
\begin{equation*}
\mathbf{B} = \mu_0(\mathbf{H}+\mathbf{M})\tag{B.23}
\end{equation*}
于是可以写出只含 $\mathbf{J}_{\text{free}}$ 的式 B.15 的新版本：不存在变化电场时
\begin{equation*}
\nabla\times\mathbf{H} = \mathbf{J}_{\text{free}}.\tag{B.24}
\end{equation*}
这很有用——它把 H 场联系到能在安培计上读出的量（$\mathbf{J}_{\text{bound}}$
可不能）。遗憾的是，这一做法简化了一条麦克斯韦方程，却复杂化了另一条：式 B.13
从简单的 $\nabla\cdot\mathbf{B}=0$ 变成略微更绕弯的
\begin{equation*}
\nabla\cdot\mathbf{H} = -\nabla\cdot\mathbf{M}.\tag{B.25}
\end{equation*}
于是 H 与 B 不同，并非无散的（用数学行话说不是螺线的）。H 的表现就好像磁单极
（即电荷的磁对应物）存在一样。

\section{物质中的麦克斯韦方程组}

存在物质时麦克斯韦方程组于是成为
\begin{align*}
\nabla\cdot\mathbf{D} &= \rho_{\mathrm{free}}\tag{B.26}\\
\nabla\cdot\mathbf{B} &= 0\tag{B.27}\\
\nabla\times\boldsymbol{\varepsilon} &= -\frac{\partial\mathbf{B}}{\partial t}\tag{B.28}\\
\nabla\times\mathbf{H} &= \mathbf{J}_{\text{free}} + \frac{\partial\mathbf{D}}{\partial t},\tag{B.29}
\end{align*}
自由电荷与束缚电荷之间作了类似的区分，电位移 D 由
\begin{equation*}
\mathbf{D} = \epsilon_0\mathbf{E} + \mathbf{P}\tag{B.30}
\end{equation*}
给出，P 是电极化强度。

\section{边界条件}

\fn{1}{在物理学文献中它们传统上被称为“药盒”（pill-boxes）——尽管药丸早已不再
用那种形状的盒子出售了！}
B 场与 H 场的一个重要差别在它们遇到两个不同区域的边界时显现。区域 1 与 2 之间的
边界条件（见图 B.4）可推导如下。考虑跨在表面上的鞋油盒状的盒子。\fn{2}{我们还要确保这一点：让盒子的厚度在半径趋于零之前先趋于零。}$\nabla\cdot\mathbf{B}=0$
结合散度定理意味着 $\int_S\mathbf{B}\cdot\mathrm{d}\mathbf{S}=0$。由于盒子的面积
主要由其顶面与底面代表，B 的垂直分量连续，即
\begin{equation*}
\mathbf{B}_{\perp 1} = \mathbf{B}_{\perp 2}.\tag{B.31}
\end{equation*}
这一论证对 $\mathbf{H}$ 不适用，因为 $\nabla\cdot\mathbf{H}=-\nabla\cdot\mathbf{M}$
可能非零：若盒内磁化强度有散度（区域 1 磁化而区域 2 未磁化时完全可能如此），
垂直分量可能不连续。平行分量的论证来自方程 $\nabla\times\mathbf{H}=\mathbf{J}_{\mathrm{free}}$，
它意味着 $\oint\mathbf{H}\cdot\mathrm{d}\mathbf{l}=\int\mathbf{J}_{\mathrm{free}}\cdot\mathrm{d}\mathbf{S}$。
当图 B.4 右侧所示的安培回路收缩到零时，回路包围的自由电流也趋于零，故 $\mathbf{H}$
的平行分量连续，即
\begin{equation*}
H_{\|1} = H_{\|2}.\tag{B.32}
\end{equation*}
这一论证对 $\mathbf{B}$ 不适用：因为 $\nabla\times\mathbf{B}=\mu_0(\mathbf{J}_{\mathrm{free}}+\mathbf{J}_{\mathrm{bound}})$，
安培回路在收缩到零的过程中虽不包围自由电流，却可能在极限下包围一些束缚面电流。

\begin{figure}
\centering
\includegraphics[width=0.45\textwidth]{figB_04.pdf}
\caption{区域 1 与 2 之间的电磁边界条件。}
\end{figure}

\section*{延伸阅读}

\begin{itemize}[itemsep=0.2em]
\item B.~I.~Bleaney 与 B.~Bleaney,《Electricity and magnetism》, OUP 1989。
\item D.~J.~Griffiths,《Introduction to electrodynamics》, Prentice-Hall 1989。
\item J.~D.~Jackson,《Classical electrodynamics》, John-Wiley 1962。
\end{itemize}
